\documentclass[11pt]{article}
\textwidth 15.9cm
\textheight 23. cm
\addtolength{\oddsidemargin}{-11.5mm}
\addtolength{\evensidemargin}{-11.5mm}
\addtolength{\topmargin}{-19.0mm}

%\usepackage{amsmath}
\usepackage[ansinew]{inputenc}
\usepackage[bbgreekl]{mathbbol}
\usepackage{amsmath,amssymb}
\usepackage{graphicx}
\usepackage{hyperref}
\usepackage{url}
\usepackage[numbers,square]{natbib}
\usepackage{multirow}
\usepackage{tikz}
\usepackage{tikz-cd}
\usetikzlibrary{arrows,matrix}
\usepackage{bbm}
\usepackage{authblk}
\usepackage{subcaption}
\renewcommand{\Affilfont}{\footnotesize}
%\renewcommand{\Authfont}{\normalsize}

%\usepackage{algpseudocode}

\usepackage{stmaryrd}
\usepackage{comment}
\usepackage{showlabels}
\usepackage{soul}
%\usepackage[dvipsnames]{xcolor}

\usetikzlibrary{calc,patterns,shapes.geometric}
\usetikzlibrary{arrows,shapes,positioning}
\usetikzlibrary{decorations.markings,decorations.pathmorphing,decorations.pathreplacing}
\usetikzlibrary{decorations.text,decorations.shapes}

%%%%%%%% MACROS
\include{gempic_spectral_macros}

\title{Hodge operators}


\author[1]{The Gempic development team}
\affil[1]{Max-Planck-Institut f\"ur Plasmaphysik, Garching, Germany}


\begin{document}

\maketitle

\begin{abstract}
This note describes some conventions that should be followed in the GEMPIC code
and its documentation.
\end{abstract}

\tableofcontents

\section{Hodge operators}
Starting with a cartesian 3D grid, a 2-form on the primal grid $\mathbf{f_2}$ such that $ \mathbf{f}^2 = \{ f_{x}^2, f_{y}^2, f_{z}^2 \}$ and a 1-form on the dual grid $\mathbf{\tilde f^1}$ such that $ \mathbf{\tilde f^1} = \{ \tilde f_{x}^1, \tilde f_{y}^1, \tilde f_{z}^1 \}$, for the calculation of the hodge operator $\mathbf{H_2}\mathbf{f^2} = \mathbf{\tilde f^1}$ we consider:
\linebreak

The NodeToCell stencils, which we call $\mathbf{w}$: 

\begin{itemize}
    \item Degree 2: $\text{stencil} = 1$
    \item Degree 4: $\text{stencil} = \{ \frac{1}{8}, \frac{6}{8}, \frac{1}{8} \}$ 
    \item Degree 6: $\text{stencil} = \{ -\frac{7}{384}, \frac{44}{384}, \frac{310}{384}, \frac{44}{384}, -\frac{7}{384} \}$ 
\end{itemize}

The CellToNode stencils, which we call $\mathbf{h}$: 

\begin{itemize}
    \item Degree 2: $\text{stencil} = 1$
    \item Degree 4: $\text{stencil} = \{ -\frac{1}{24}, \frac{13}{12}, -\frac{1}{24} \}$ 
    \item Degree 6: $\text{stencil} = \{ \frac{3}{640}, -\frac{29}{480}, \frac{1067}{960}, -\frac{29}{480}, \frac{3}{640} \}$ 
\end{itemize}

The stencils will have subindices $i,j,k$ depending on whether they are applied to the $x, y, z$ direction. The degree is fixed from the beginning.
\linebreak

Therefore, the hodge formula for this map is the following for the x component, considering that for degree 2, 4 and 6, the values for $d$ are $0, 1$ and $2$ respectively:

\begin{equation}
    \sum_{s_i = -d} ^ {+d} \sum_{s_j = -d} ^ {+d} \sum_{s_k = -d} ^ {+d} w(s_i)h(s_j)h(s_k) f_{x}^2 (i + s_i, j + s_j, k + s_k) = \frac{\Delta y \Delta z}{\Delta x} f_{x}^1 (i, j, k)
\end{equation}

For the y component:

\begin{equation}
    \sum_{s_i = -d} ^ {+d} \sum_{s_j = -d} ^ {+d} \sum_{s_k = -d} ^ {+d} h(s_i)w(s_j)h(s_k) f_{y}^2 (i + s_i, j + s_j, k + s_k) = \frac{\Delta x \Delta z}{\Delta y} f_{y}^1 (i, j, k)
\end{equation}

For the z component:

\begin{equation}
    \sum_{s_i = -d} ^ {+d} \sum_{s_j = -d} ^ {+d} \sum_{s_k = -d} ^ {+d} h(s_i)h(s_j)w(s_k) f_{z}^2 (i + s_i, j + s_j, k + s_k) = \frac{\Delta x \Delta y}{\Delta z} f_{z}^1 (i, j, k)
\end{equation}

\subsection{Hodge algorithm}

We transform the expression (mind we use square brackets to denote an array of indices for the form)
\begin{equation}
    \sum_{s_i = -d} ^ {+d} \sum_{s_j = -d} ^ {+d} \sum_{s_k = -d} ^ {+d} w[s_i]h[s_j]h[s_k] f_{x}^2 [i + s_i, j + s_j, k + s_k]
\end{equation}

into

\begin{equation}
    t_1 [i,j,k] = \sum_{s_k} h[s_k] f_x ^ 2[i, j, k + s_k]
\end{equation}

Since the sums are independent on each other:
\begin{equation}
    t_2 [i,j,k] = \sum_{s_j} h[s_j] t_1[i, j + s_j, k]
\end{equation}

\begin{equation}
    t_3 [i,j,k] = \sum_{s_i} w[s_i] t_2[i + s_i, j, k]
\end{equation}

Where $t_1$ is a reference to the output, $t_2$ is a temporary MultiFab and $t_3$ is also a reference to the output. This avoids declaring three temporary MultiFabs (which is potentially very expensive and inefficient).


\end{document}
